% =====================================================================
%  §1  Abstract and positioning
% =====================================================================
\section{Abstract and positioning: the continuum-limit coordinate is a consequence, not an entrance}
\label{sec:position}

\begin{keybox}
\textbf{One sentence.} \emph{All ``shape'' and ``scaling'' arguments are logically posterior to
$m>0$; therefore the continuum-limit coordinate is not where the gap is proved.} \LN
\end{keybox}
\src{\code{LOOP-RUNBOOK.md} \S10.21(5) (verbatim: ``any scaling statement about $m$ is logically
posterior to $m>0$''), \S10.22 (``the continuum limit is not where one proves the gap; it is where
the gap's consequences live''); immune family \textbf{47} (\code{LOOP-RUNBOOK.md} line 1126); the
third admission question for a new species (ibid.\ line 1141); \code{H1B-FIVE.md} \S4.3. The corpus
tag for this verdict is \code{[NG]}.}

\begin{keybox}
\textbf{Scope, which no section below retracts.} These sections are about $\Gpp$ and about the coordinate
$m/\Lm$; \emph{they do not imply a mass gap.} Two \LT{}-level obstructions are recorded in
\code{docs/paper-v5} \S2--\S3: (i) $\Gpp=\Delta_{\mathrm{dyn}}+\Wf$ with $\Wf\ge N/12>0$, so $\Gpp>0$ is
strictly weaker than $\Delta_{\mathrm{dyn}}>0$; (ii) the harmonic-oscillator~$\otimes$~free-mode
counterexample has $\Gpp=2\omega>0$ while $m=0$. The correct target is $f_{1}^{2}$ (the
$O=\frac14\operatorname{tr}F^{2}$ pole residue) together with $(CP)$; see \code{paper-v5} \S7. \LN
\end{keybox}
\src{\code{docs/paper-v5} \S2--\S3 and \S7; \code{T01-IMPLY.md}, \code{T03-TOY.md}. Every label in the
sections below stands as stated; withdrawn is only the reading of $\Gpp>0$ as a route to $m>0$.}

\subsection{What is refuted is the ``route'' reading, not the shape}

On 2026-11-25 this campaign fixed its single objective as
\begin{equation}
  \frac{m}{\Lm} \;\ge\; c(N)\;>\;0 ,
  \qquad c \ \text{a pure number},
  \label{eq:target}
\end{equation}
and reduced it to a slope identity along $\beta$,
$\dd\ln(ma)/\dd\beta = -12\pi^{2}/(11N^{2})$ (\code{CONTINUUM-PLAN.md} \S1). The reduction itself is
an identity: $1/(2b_0g^{2}) = 12\pi^{2}\beta/(11N^{2})$ with $b_0=11N/(48\pi^{2})$. \LT
\src{\code{CONTINUUM-PLAN.md} \S1 (including the $\mathrm{SU}(2)$ coefficient $3\pi^{2}/11 = 2.691710$
and the $\mathrm{SU}(3)$ coefficient $4\pi^{2}/33 = 1.196316$); independently rechecked in
\code{LOOP-RUNBOOK.md} \S10.19 and \S10.21(1).}

But reading it as a \emph{route} (``prove the slope identity $\Rightarrow$ obtain the gap'') died the
same day, for three independent reasons. \LN
\src{\code{C1-SLOPE-EQUIV.md}, as reported in \code{LOOP-RUNBOOK.md} \S10.21; the three reasons are
listed below one by one.}

\begin{enumerate}[label=(\Alph*)]
\item \textbf{Same order $=$ tautology.} Write $\hat m := m a$. Since
      $\hat m = (m/\Lambda)(a\Lambda)$, one has $\hat m/(a\Lambda) \equiv m/\Lambda$ ---
      \emph{zero physical content}. \LT
\item \textbf{Mixed order $=$ false, in both directions.} A perfectly physical $m/\Lambda = 1$ gives
      the slope $-1.126068 \neq -1.196316$; conversely, forcing the slope to be the constant forces
      $m/\Lambda \propto \beta^{-0.4215}$, which is false as well. \LT
\item \textbf{Presupposition.} The slope identity presupposes $ma>0$; every scaling statement about
      $m$ is logically posterior to $m>0$, hence it cannot structurally be used to prove a gap. \LN
\end{enumerate}
\src{\code{LOOP-RUNBOOK.md} \S10.21(2)(3)(4)(5); the numbers $1.126068$ and $0.4215$ are the two
counterexamples of item (4) there.}

\subsection{What survives is uniqueness: the shape is a theorem, and a theorem is not a route}

The route reading is dead; the \emph{shape} is promoted to a theorem. The list of scale sources is
closed: $\Lambda$ is the only independent scale ($a$ is only a unit and enters only through
$\Lambda a$; $L$ enters only through $\Lambda L$; $\theta$ is an angle with $[\,\theta\,]=0$; $N$ is
discrete; $T$ disappears at $T=0$), and $\theta=0$ holds \emph{by definition} from the positivity of
the Wilson measure, not by hypothesis. Hence
\begin{equation}
  m \;=\; C(N,\ \text{scheme})\cdot\Lambda
  \qquad\text{is the only possible shape},
  \label{eq:unique}
\end{equation}
escape shapes $m=C/a$ and $m=C\Lambda^{2}/g$ are excluded, and $C$ may depend on $N$, on $\theta=0$
and on the scheme ($\Lm/\Lambda_I = 2.633$) but \emph{not} on $\beta,a,L,T$. \LT
\src{\code{C3-UNIQUENESS.md}, as reported in \code{LOOP-RUNBOOK.md} \S10.19 and in the first row of
the table in \S10.22; the scheme ratio $2.633$ is from \code{L2-GLUEBALL-LAMBDA.md}.}

\noindent
This is exactly what the title means: \eqref{eq:unique} is a theorem about how $m$ must look
\emph{given that $m$ exists}. \textbf{A theorem about shape can never serve as the route that
produces $m>0$} --- it is a consequence. \LT
\src{Ibid.; see also the third admission question in \code{LOOP-RUNBOOK.md} line 1141 (``does it
deliver $m>0$ or only some property of $m$?'').}

\subsection{The same day turned ``no work inside the continuum limit'' into five closures}

The remaining sections are not new conjectures; they cash out the verdict coordinate by coordinate.
The last row of the table is the only thing still alive.

\begin{table}[h]
\centering\footnotesize
\setlength{\tabcolsep}{4pt}
\renewcommand{\arraystretch}{1.15}
\begin{tabularx}{\textwidth}{@{}l l X l@{}}
\toprule
\textbf{\S} & \textbf{Object} & \textbf{Verdict (one line)} & \textbf{Label}\\
\midrule
\S\ref{sec:scalar} & scalar form of H1 & $\Var - E[H] - \Wf = -\Gpp$, hence $H1(\delta) \iff \Gpp \ge \delta$: there is no independent variance inequality & \LT\\
\S\ref{sec:failure} & solvable models & $\delta(\beta)=0.2051153-0.0110593\beta^{2}$, $\beta^{*}=4.30661$, but at $N=3$ the coefficient changes sign, so H1 \textbf{holds on the whole $\mathrm{SU}(3)$ scaling region}; with $\beta^{*}(N)\propto N^{3/2}$ one gets $\beta^{*}(3)\approx19.6\gg\beta_c(\mathrm{SU}(3))=5.6924$ and $\beta^{*}/\beta_c=1.87\to3.45$, so the exact failure sits outside the scaling region of \emph{both} groups, and the $\beta$-uniform reading is dead & \LE{}/\LN\\
\S\ref{sec:five} & five conditions & $C5\Rightarrow C2$; $C3/C4$ are filters; only $C5$'s rate uniformisation $=d5/(CP)=p=1$ survives & \LM\\
\S\ref{sec:tension} & string-tension side & $\sigma\propto 1/a^{2}$ diverges as a power; $\sigma=2g^{2}N m^{2}$ is scheme bookkeeping; $\sigma>0\not\Rightarrow m>0$ & \LN\\
\S\ref{sec:loop} & loop algebra & $\kappa_W=0$ excludes homogeneous monomials only (a weak result); the area law holds only in the degenerate sector; the non-degenerate sector is never elliptic & \LN\\
\S\ref{sec:lit} & literature & $m/\Lm = 6.3(3)$ at the continuum limit, three routes agreeing; $a$-uniform theorems: \textbf{zero} & \LE\\
\S\ref{sec:errata} & errata & 24 normalization / coordinate / convention errors of the day, each as ``wrong $\to$ right $\to$ source'' & ---\\
\bottomrule
\end{tabularx}
\caption{The structure of this document. The first five sections are closures, the sixth is the only
surviving target, the seventh is the external calibration, the eighth is the ledger.}
\label{tab:map}
\end{table}

\subsection{Conventions (read before \S\ref{sec:tension} and \S\ref{sec:errata})}

\begin{itemize}
\item Lattice $\beta := 2N/g^{2}$ (Wilson convention); $\tH := \beta N = 2N^{2}/g^{2}$.
\item SZZ's $\beta_A$ is \textbf{not} the same coordinate: $\beta_W = N^{2}\beta_A$. Mixing them
      misses by exactly $N^{2}$ ($\mathrm{SU}(3)$: $9$); this is errors 4 and 10 of the day. \LT
      \src{\code{BETA-CONVENTION.md} (read to arXiv:2204.12737v1 \S1.1 eq.~(1.2), \S1.2 eq.~(1.3));
      \code{LOOP-RUNBOOK.md} \S9.8.}
\item Every $\delta$ below is a \textbf{surfaced lower bound in lattice units} (dimensionless) and
      must not be read as a physical mass bound: $\Gpp_a = a^{2}m_{\mathrm{phys}}^{2}$, so ``for
      fixed $a$, $\Gpp_a \ge \lambda$'' is dimensionally wrong; the correct target is
      $\Gpp_a/a^{2} \ge \Delta_0^{2}$. \LT
      \src{\code{docs/paper-v2/ch03-existence.tex}:146--165 (\code{sec:dim}), as reported in
      \code{H1A-EXACT.md} \S1.4.}
\end{itemize}
